\documentclass[10pt,a4paper]{article}
% Based on the supplied Machine Learning Mini-Project LaTeX template.
% Compile from this folder: pdflatex SCF_MiniProject_Report.tex
\usepackage[margin=20mm]{geometry}
\usepackage[T1]{fontenc}
\usepackage{newtxtext}
\usepackage{newtxmath}
\usepackage{microtype}
\usepackage{graphicx}
\usepackage{booktabs}
\usepackage{amsmath}
\usepackage[colorlinks=true,urlcolor=blue!55!black,linkcolor=blue!55!black]{hyperref}
\usepackage{xcolor}
\setlength{\parindent}{0pt}
\setlength{\parskip}{3pt}
\setlength{\textfloatsep}{7pt}
\setlength{\intextsep}{7pt}
\renewcommand{\refname}{References}

\begin{document}
\begin{center}
{\Large\bfseries Predicting Stress Concentration in CFST T-Joints with Regression\par}
\vspace{5pt}
{\normalsize Ye Zining\par}
\vspace{2pt}
{\small Course: Machine Learning Mini-Project \quad | \quad GitHub: \href{https://github.com/yejacob18-prog/CFST-SCF-ML-MiniProject}{CFST-SCF-ML-MiniProject}\par}
\end{center}
\vspace{-6pt}

\section{Problem and Motivation}
Stress is not distributed evenly around a welded tubular joint. When a brace transfers force into the main tube (the chord), some locations carry much higher stress than others. Engineers describe this increase using the stress concentration factor (SCF), the ratio of local hot-spot stress to nominal stress. In this project, I wanted to see whether a relatively simple machine learning model could predict the \emph{brace-saddle} SCF of a concrete-filled steel tubular (CFST) T-joint under brace compression. The inputs are three geometric ratios, $\beta$, $\gamma$ and $\tau$, and the output is a numerical SCF. If the predictions are reasonably accurate, such a model could help with quick comparisons of joint geometries before more detailed analysis.

\section{Dataset and Method}
The data came from a publicly released CFST joint dataset by Bajracharya \emph{et al.}~\cite{dataset,paper}. I used the sheet \texttt{T-JOINT BS IN COMP}, which contains \textbf{140 finite-element-analysis (FEA) results} covering a $4\times5\times7$ grid of geometric ratios: $\beta=0.3$--$0.6$, $\gamma=20$--$40$ and $\tau=0.4$--$1.0$. The FEA column was used as the target SCF (1.60--2.14). I did not use the original authors' ANN predictions as inputs, since that would leak prediction information into the new model. There were no missing values or duplicate geometric combinations.

I randomly set aside \textbf{20\% of the data (28 cases) for final testing}, using seed 123. The other 112 cases were used for training and \textbf{five-fold cross-validation}. I compared a mean-value baseline, linear ordinary least squares (OLS), linear stochastic gradient descent (SGD), and second- or third-degree polynomial models fitted with OLS, Ridge ($L_2$) or Lasso ($L_1$). Cross-validation was used to select the polynomial degree, regularization strength and SGD learning rate. Feature expansion and standardization were done inside the training folds to avoid data leakage. After model selection, I refitted the chosen model on all 112 training cases and tested it on the untouched 28 cases. Performance was measured using mean absolute error (MAE), root-mean-square error (RMSE) and $R^2$. The code uses scikit-learn~\cite{sklearn}.

\section{Results and Discussion}
\begin{table}[ht]
\centering
\caption{Five-fold results from the training set. Lower MAE means better predictions.}
\small
\begin{tabular}{lrrl}
\toprule
Model & Train MAE & Val. MAE & Selected setting\\
\midrule
Mean baseline & 0.10615 & 0.10819 & Mean\\
Linear OLS & 0.03525 & 0.03607 & Degree 1\\
Linear SGD & 0.03525 & 0.03608 & $\eta_0=0.01$\\
Polynomial OLS & 0.01160 & \textbf{0.01722} & Degree 3\\
Polynomial Ridge & 0.01185 & 0.01723 & Degree 3, $\alpha=0.001$\\
Polynomial Lasso & 0.01545 & 0.01853 & Degree 3, $\alpha=0.0001$\\
\bottomrule
\end{tabular}
\end{table}

The linear models were already much better than predicting the average SCF, but adding polynomial terms reduced validation MAE further, from 0.03607 to 0.01722. The best result came from a third-degree polynomial OLS model. Ridge was almost identical (0.01723), while Lasso was slightly worse (0.01853). On the \textbf{28 unseen test cases}, the chosen model reached \textbf{MAE = 0.01973, RMSE = 0.02755 and $R^2=0.96693$}. Most predictions lie close to the ideal line in Fig.~\ref{fig:test}, although not every case was predicted equally well. For example, at $(\beta,\gamma,\tau)=(0.6,40,0.4)$, the model predicted 1.78 against an FEA value of 1.87, an error of about 0.09. The gap between training and validation MAE also suggests that the more flexible fit needs some caution.

\clearpage
\begin{figure}[ht]
\centering
\includegraphics[width=0.46\linewidth]{../figures/test_prediction_scatter.png}
\caption{Predicted versus FEA SCF for the 28 held-out cases. The dashed line shows perfect agreement.}
\label{fig:test}
\end{figure}

\section{Conclusions and Reflection}
Choosing the right model complexity is more important than using a more complicated training method. Linear regression using ordinary least squares (OLS) is fast and easy to understand, but it cannot represent nonlinear relationships well. In this project, stochastic gradient descent (SGD) gave almost the same results as linear OLS, suggesting that iterative optimization offered little additional benefit for this small dataset. The cubic polynomial regression model was more flexible and provided much better predictions. However, adding more polynomial terms also increases the risk of overfitting. Ridge regression and Lasso regression can reduce this risk by shrinking the regression coefficients through regularization, but they require additional hyperparameter tuning. In this experiment, neither method showed a clear improvement over polynomial regression using OLS.

In this project, the final model can be used for quick estimations within the existing range of geometric parameters, but it is not yet sufficiently validated for use in new joint designs. This is mainly because the available data is limited. All 140 data points were obtained from finite element analysis (FEA) using a predefined grid of geometric parameter combinations, and the test set contained only 28 randomly held-out combinations. Therefore, the predictive accuracy of the current model for new geometries outside the existing parameter grid remains uncertain, and the model has not been validated against experimental measurements. One possible improvement would be to generate more geometric configurations with new combinations of geometric parameters and perform additional FEA simulations. If possible, experimental testing of joint specimens could also be conducted to obtain actual measurement data. The FEA results, model predictions, and experimental measurements could then be compared to identify the differences between them. This would provide a more realistic evaluation of the model's predictive accuracy and generalization ability, helping us better understand under which conditions the model can be considered reliable.

\section*{Reproducibility and Disclosures}
The linked GitHub repository provides the cleaned CSV, training and prediction scripts, model settings, evaluation results and README instructions. The original workbook is available from the cited Mendeley Data record. No pretrained model or original ANN code was reused. ChatGPT assisted with Python code, and the numerical results were obtained by running the project code and could be checked.

\begingroup\small\begin{thebibliography}{9}\setlength{\itemsep}{0pt}
\bibitem{dataset} S. Bajracharya, S. Nakamura, and T. Nishikawa. \emph{Dataset for SCF Prediction in CFST T- and K-Joints Using ANN}. Mendeley Data, V1 (2025). \url{https://doi.org/10.17632/rp7vm5627z.1}.
\bibitem{paper} S. Bajracharya, S. Nakamura, and T. Nishikawa. \emph{Prediction of stress concentration factors for concrete-filled steel tubular (CFST) T-joints and K-joints using artificial neural network}. Results in Engineering 27, 106313 (2025). \url{https://doi.org/10.1016/j.rineng.2025.106313}.
\bibitem{sklearn} F. Pedregosa \emph{et al.} Scikit-learn: Machine Learning in Python. \emph{JMLR} 12, 2825--2830 (2011).
\end{thebibliography}\endgroup
\end{document}